\begin{abstract}
Expert demonstrations are the one input to imitation learning that compute
cannot make cheaper. Once their number is fixed, what each one contains decides
the result. Existing methods answer one question well: when to ask the expert.
They call the expert at the first state where the policy is uncertain. Two
further questions are left at their defaults: which failure to correct, and
where the demonstration begins. We present \textbf{DISEIL}
(\textbf{D}emonstration d\textbf{I}stillation for \textbf{S}ample-\textbf{E}fficient
\textbf{I}mitation \textbf{L}earning). DISEIL keeps the uncertainty test and
answers both. A training round ends with failed episodes, runs that end without
achieving the task. DISEIL sorts them into modes. For each mode it prescribes
the start configuration of a demonstration not yet collected. It checks that
the configuration is feasible before any expert time is spent. We evaluate
across five tasks, with image and state observations, at a budget of 20
demonstrations, against the best baselines. That gives ten settings. DISEIL
reaches the highest mean success rate in all ten, on average about 3 percentage
points above the strongest baseline.
\end{abstract}
